% The fate of the nine elements of 𝔽₉ = ℤ[i]/(3) under the Frobenius
% x ↦ x³: 0, 1, 2 are fixed and the other six swap in pairs.
% Source: book/tex/alg-NT/finite-field.tex (§ The Galois theory of
% finite fields) — tikzcd.
\documentclass[border=2mm]{standalone}
\usepackage{tikz}
\usetikzlibrary{arrows.meta}
\begin{document}
\begin{tikzpicture}[>=latex]
  \node (z) at (0,1.4) {$0$};
  \node (o) at (1.2,1.4) {$1$};
  \node (t) at (2.4,1.4) {$2$};
  \node (i1) at (4.0,1.4) {$i$};
  \node (i2) at (4.0,0) {$-i$};
  \node (j1) at (5.8,1.4) {$1+i$};
  \node (j2) at (5.8,0) {$1-i$};
  \node (k1) at (7.8,1.4) {$2+i$};
  \node (k2) at (7.8,0) {$2-i$};
  \draw[<->] (i1) -- (i2) node[midway,right] {$\sigma$};
  \draw[<->] (j1) -- (j2) node[midway,right] {$\sigma$};
  \draw[<->] (k1) -- (k2) node[midway,right] {$\sigma$};
\end{tikzpicture}
\end{document}
